\documentclass[11pt]{article}
\usepackage[margin=1in]{geometry}
\usepackage{enumitem}
% =============================================================================
% Preambles/header.tex — shared LaTeX/Beamer preamble for this project
%
% Use from a Beamer lecture by prepending:
%     \input{header}
% and compiling with TEXINPUTS=../Preambles:$TEXINPUTS (the /compile-latex
% skill does this automatically).
%
% PALETTE CONTRACT. Color names defined here MUST match the names in
% Quarto/theme-template.scss so Beamer and Quarto renderings use the same
% palette. The scripts/check-palette-sync.sh script enforces this.
%
% To customize: edit the HEX values below AND the matching $variable in
% Quarto/theme-template.scss, then run scripts/check-palette-sync.sh.
% =============================================================================

% -----------------------------------------------------------------------------
% Engine detection + encoding
%
% Our /compile-latex skill uses XeLaTeX, where inputenc is unnecessary and
% can produce encoding warnings. Only load inputenc under pdfTeX.
% -----------------------------------------------------------------------------
\usepackage{iftex}
\ifPDFTeX
  \usepackage[utf8]{inputenc}
\fi

\usepackage{amsmath, amssymb, amsthm}
\usepackage{booktabs}
\usepackage{graphicx}

% -----------------------------------------------------------------------------
% PALETTE — mirrors Quarto/theme-template.scss variable names.
% Load xcolor and define colors BEFORE anything that references them
% (notably hyperref, hypersetup, and the Beamer color assignments).
% -----------------------------------------------------------------------------
\usepackage{xcolor}

\definecolor{primary-blue}{HTML}{012169}      % headings, accents
\definecolor{primary-gold}{HTML}{B9975B}      % emphasis, borders
\definecolor{highlight-yellow}{HTML}{F2A900}  % markers, alerts
\definecolor{light-bg}{HTML}{E8EDF5}          % title / section backgrounds
\definecolor{jet}{HTML}{1A1A1A}               % body text

% Semantic colors (used in diagrams and annotations)
\definecolor{positive}{HTML}{15803D}          % good / observed / identified
\definecolor{negative}{HTML}{B91C1C}          % bad / problematic / bias
\definecolor{neutral}{HTML}{525252}           % reference / context

% Highlight accents (match .hi-* classes in the SCSS)
\definecolor{hi-slate}{HTML}{314F4F}
\definecolor{hi-green}{HTML}{2E8B57}
\definecolor{hi-red}{HTML}{C41E3A}

% Semantic aliases so skill templates and snippets can write
% \textcolor{accent}{...} without hard-coding "primary-blue".
\colorlet{accent}{primary-blue}
\colorlet{accent2}{primary-gold}

% -----------------------------------------------------------------------------
% Hyperref — loaded AFTER xcolor + palette so link colors resolve.
% -----------------------------------------------------------------------------
\usepackage{hyperref}
\hypersetup{colorlinks=true, linkcolor=primary-blue, urlcolor=primary-blue,
            citecolor=primary-blue, pdfborder={0 0 0}}

% -----------------------------------------------------------------------------
% Course code — set once per deck (after \input{header}, before \begin{document})
% via \coursecode{CS401}. Shown in the footer on every slide so a deck's course
% is identifiable without opening the title page. Empty by default.
% -----------------------------------------------------------------------------
\makeatletter
\newcommand{\coursecode}[1]{\gdef\@coursecodetext{#1}}
\gdef\@coursecodetext{}
\makeatother

% -----------------------------------------------------------------------------
% Beamer theme assignments (only apply under Beamer).
%
% We detect Beamer via \@ifclassloaded{beamer}, which is the documented,
% non-internal way. This must be wrapped in \makeatletter/\makeatother
% because \@ifclassloaded contains an @-catcode character.
% -----------------------------------------------------------------------------
\makeatletter
\@ifclassloaded{beamer}{%
  \usefonttheme{professionalfonts}
  \setbeamercolor{structure}{fg=primary-blue}
  \setbeamercolor{frametitle}{fg=primary-blue}
  \setbeamercolor{title}{fg=primary-blue}
  \setbeamercolor{subtitle}{fg=primary-gold}
  \setbeamercolor{author}{fg=primary-blue}
  \setbeamercolor{institute}{fg=neutral}
  \setbeamercolor{date}{fg=primary-gold}
  \setbeamercolor{itemize item}{fg=primary-blue}
  \setbeamercolor{itemize subitem}{fg=primary-gold}
  \setbeamercolor{alerted text}{fg=primary-gold}
  \setbeamercolor{block title}{fg=white, bg=primary-blue}
  \setbeamercolor{block body}{bg=light-bg}
  % Semantic block variants: block=definition/concept (blue), exampleblock=
  % worked example (green/positive), alertblock=key takeaway or caution (gold).
  % Keep to <=2 colored boxes per slide across ALL three variants (INV-7).
  \setbeamercolor{block title example}{fg=white, bg=positive}
  \setbeamercolor{block body example}{bg=light-bg}
  \setbeamercolor{block title alerted}{fg=white, bg=primary-gold}
  \setbeamercolor{block body alerted}{bg=light-bg}

  % Minimal footer: course code (left) + page X / N (right), no navigation clutter
  \setbeamertemplate{navigation symbols}{}
  \setbeamertemplate{footline}{%
    \color{neutral}\footnotesize\hspace{0.7em}\@coursecodetext\hfill
    \insertframenumber/\inserttotalframenumber\hspace{0.7em}\vspace{0.3em}}
}{}
\makeatother

% -----------------------------------------------------------------------------
% TikZ setup — libraries the snippet gallery uses
% -----------------------------------------------------------------------------
\usepackage{tikz}
\usetikzlibrary{arrows.meta, positioning, calc, decorations.pathreplacing,
                fit, shapes.geometric, backgrounds}

% Shared TikZ styles — snippets and hand-written diagrams can pick these up
% via `\begin{tikzpicture}[every node/.style={...}]` or by naming specific
% styles (e.g., `\node[dag-node] (X) at (0,0) {X};`).
\tikzset{
  % Node styles for causal DAGs and similar
  dag-node/.style = {
    draw=primary-blue, fill=light-bg, rounded corners=2pt,
    minimum width=1.2cm, minimum height=0.9cm, align=center, font=\small
  },
  decision-node/.style = {
    draw=primary-gold, fill=white, diamond, aspect=1.8,
    minimum width=1.8cm, minimum height=1.2cm, align=center, font=\small,
    inner sep=1pt
  },
  % Edge styles — observed vs counterfactual
  observed-edge/.style = {-{Latex[length=2.5mm]}, thick, draw=jet},
  counterfactual-edge/.style = {-{Latex[length=2.5mm]}, thick, draw=neutral, dashed},
  confound-edge/.style = {-{Latex[length=2.5mm]}, thick, draw=negative, dashed},
  % Data-point styles for plots
  observed-dot/.style = {fill=primary-blue, draw=primary-blue, circle, inner sep=1.2pt},
  counterfactual-dot/.style = {fill=white, draw=neutral, circle, inner sep=1.2pt}
}

% -----------------------------------------------------------------------------
% Convenience macros
% -----------------------------------------------------------------------------
\newcommand{\muted}[1]{\textcolor{neutral}{#1}}
\newcommand{\key}[1]{\textcolor{primary-gold}{\textbf{#1}}}
\newcommand{\good}[1]{\textcolor{positive}{#1}}
\newcommand{\bad}[1]{\textcolor{negative}{#1}}

% Standout frame macro: full-bleed dark background for section transitions.
% Beamer-only (uses \begin{frame}/\setbeamercolor/\usebeamerfont) -- do not
% call from an article-class file (e.g. Notes/<CODE>/*-notes.tex). \muted,
% \key, \good, \bad, and \coursecode above are plain xcolor/gdef and are
% safe to use from either class; verified when Notes/article-class support
% was added.
% Expand: \transitionslide{Your transition title}
\newcommand{\transitionslide}[1]{%
  \begingroup
  \setbeamercolor{background canvas}{bg=primary-blue}
  \begin{frame}[plain]
    \centering
    \vfill
    {\usebeamerfont{title}\color{white}\Large #1\par}
    \vfill
  \end{frame}
  \endgroup
}

% Section divider frame (Beamer-only), an alternative to \transitionslide with a
% darker two-line style: near-black (jet) background, a small white label line
% above a larger gold title, and a thin gold rule along the bottom edge.
% Expand: \sectiondivider{Section 2}{Pointers}
\newcommand{\sectiondivider}[2]{%
  \begingroup
  \setbeamercolor{background canvas}{bg=jet}
  \begin{frame}[plain]
    \vspace*{3.0cm}
    {\color{white}\ttfamily\large #1\par}
    \vspace{0.5cm}
    {\color{primary-gold}\LARGE #2\par}
    \begin{tikzpicture}[remember picture, overlay]
      \draw[primary-gold, line width=2.5pt]
        ($(current page.south west)+(0,0.1)$) -- ($(current page.south east)+(0,0.1)$);
    \end{tikzpicture}
  \end{frame}
  \endgroup
}

% -----------------------------------------------------------------------------
% Channel watermark — "Concepts That Click" (YouTube).
%
% Article-class files (CompetitiveExam Q/A sets, Notes, Assignments) get a
% light diagonal draft watermark on every page plus a centered footer naming
% the channel. Beamer decks get a subtle TikZ background watermark instead
% (the footline already carries the course code). Centralized here so every
% MCQ PDF is branded without per-file edits.
% -----------------------------------------------------------------------------
\makeatletter
\@ifclassloaded{article}{%
  \usepackage{draftwatermark}
  \SetWatermarkText{Concepts That Click}
  \SetWatermarkScale{1}
  \SetWatermarkFontSize{2cm}
  \SetWatermarkLightness{0.86}
  \SetWatermarkAngle{45}
  \usepackage{fancyhdr}
  \pagestyle{fancy}
  \fancyhf{}
  \fancyfoot[C]{\footnotesize\textcolor{neutral}{Concepts That Click~|~YouTube: \href{https://www.youtube.com/@concepts.thatclick}{@concepts.thatclick}}}
  \renewcommand{\headrulewidth}{0pt}
  \renewcommand{\footrulewidth}{0pt}
  \fancypagestyle{plain}{%
    \fancyhf{}%
    \fancyfoot[C]{\footnotesize\textcolor{neutral}{Concepts That Click~|~YouTube: \href{https://www.youtube.com/@concepts.thatclick}{@concepts.thatclick}}}%
    \renewcommand{\headrulewidth}{0pt}%
    \renewcommand{\footrulewidth}{0pt}%
  }
}{}
\@ifclassloaded{beamer}{%
  \setbeamertemplate{background}{%
    \begin{tikzpicture}[remember picture, overlay]
      \node[rotate=45, opacity=0.055, align=center]
        at (current page.center)
        {\Huge\textcolor{neutral}{Concepts That Click}};
    \end{tikzpicture}%
  }
}{}
\makeatother 
\coursecode{JKSSB}
\renewcommand{\thesection}{44.\arabic{section}}
\newtheorem{example}{Example}[section]
\newenvironment{definitionbox}[1]{%
  \par\noindent\textbf{Definition (#1).}\ \itshape
}{\par\vspace{0.3em}}
\title{Cybersecurity Prevention --- Detailed Lecture Notes}
\author{JKSSB Computer Awareness}
\date{\today}

\begin{document}
\maketitle

\section{Prevention and the CIA triad}
Cleaning an infected machine is slow and uncertain, so the first goal of
security is prevention: small repeated habits plus a few tools that each do
one distinct job. The standard way to judge whether protection is complete is
the \textbf{CIA triad}: \textbf{confidentiality}, \textbf{integrity}, and
\textbf{availability}.

\begin{definitionbox}{Confidentiality}
Only authorised persons can read the data. Passwords, encryption, and MFA
protect it.
\end{definitionbox}

\begin{definitionbox}{Integrity}
The data is accurate and has not been changed by unauthorised means.
Clean backups, digital signatures, and timely patching protect it.
\end{definitionbox}

\begin{definitionbox}{Availability}
The data and systems are reachable when needed. Backups, antivirus software,
and firewalls protect it.
\end{definitionbox}

\begin{center}
\begin{tabular}{p{0.24\textwidth}p{0.66\textwidth}}
\toprule
\textbf{Term} & \textbf{Meaning in this lesson} \\
\midrule
Antivirus & Software that detects, blocks, and removes malicious programs. \\
Signature & A stored pattern of a known virus, used for recognition. \\
Real-time scanning & Checking files at the moment they are opened, downloaded, or copied. \\
Quarantine & Holding a suspicious file in an isolated area so it cannot run. \\
Firewall & A filter that allows or blocks network traffic by rule. \\
Patch & A small vendor fix for a known security weakness. \\
MFA --- Multi-Factor Authentication & Proving identity with two or more independent factors. \\
Backup & A separate clean copy of data kept for recovery. \\
HTTPS & The secure, encrypted form of web browsing. \\
WPA2/WPA3 & Encryption standards that protect a Wi-Fi link. \\
Autorun & The feature that launches programs automatically when a drive is connected. \\
\bottomrule
\end{tabular}
\end{center}

One attack can break more than one leg of the triad. Stealing a password over
open Wi-Fi breaks confidentiality; altering marks in a result file breaks
integrity; ransomware that locks every document breaks availability. The rest
of this lesson maps each preventive control to the goal it serves.

\section{Antivirus: signatures, scanning, and quarantine}
\textbf{Antivirus} software detects, blocks, and removes malicious programs
(malware). Its oldest method is the \textbf{signature}: a stored pattern of a
known virus, like a fingerprint kept on record. The antivirus compares each
file against its signature database and flags anything that matches. Because
new malware appears every day, the signature database must be updated
regularly; an outdated antivirus misses new threats.

\begin{definitionbox}{Signature}
A stored pattern identifying a known malicious program, against which files
are compared.
\end{definitionbox}

Signatures catch known malware, but they cannot recognise a brand-new
(zero-day) threat. Behaviour-based detection watches what a program
\emph{does} --- for example, trying to encrypt every document --- and stops
it even when no signature matches.

Two scanning modes matter. \textbf{Real-time scanning} checks files at the
moment they are opened, downloaded, or copied, before harm is done. A
\textbf{full (on-demand) scan} checks the whole system when the user asks for
it. When either mode finds something suspicious, the antivirus moves the file
to \textbf{quarantine}: an isolated area where the file cannot run or spread.
Quarantine is not deletion. The file is held safely until the user deletes it
or restores it (for example, when a clean file was flagged by mistake).

\begin{example}
A pen drive is plugged into the office desktop and a file on it matches a
known virus signature. Real-time scanning blocks the file the moment it is
opened and moves it to quarantine, so it cannot run. The user then deletes the
quarantined file. If instead the file had been a clean office template flagged
by mistake, quarantine would allow restoring it --- deletion would not.
\end{example}

\section{Firewall: the network gatekeeper}
A \textbf{firewall} monitors incoming and outgoing network traffic and blocks
whatever its rules forbid. Its job is to stop unauthorised network access: a
stranger on the internet reaching into the office computer, or a malicious
program on the computer sending data out. A firewall can be software running
on the computer or hardware placed at the network boundary; offices commonly
use both.

\begin{definitionbox}{Firewall}
A hardware or software filter that allows or blocks network traffic
according to a set of rules.
\end{definitionbox}

The firewall--antivirus boundary is a favourite exam trap. A firewall blocks
bad \emph{traffic}; it does not remove viruses already saved on the disk.
Antivirus removes bad \emph{files}; it does not, by itself, patrol the network
boundary. The two complement each other and neither replaces the other.

\section{Patching and updates}
A \textbf{patch} is a small fix released by the vendor for a known security
weakness, and a \textbf{software update} delivers such fixes (plus features).
Attackers study published weaknesses, so any delay in installing patches
leaves known doors open. The operating system, the web browser, and the
antivirus itself must all be updated; each has flaws attackers can exploit.
Automatic updates are the safest setting for most users.

\begin{example}
The office desktop offers an operating-system update. Installing it closes a
published weakness the same week it became known. Postponing it for months
means every attacker who read the same notice knows exactly which door is
still open on that machine.
\end{example}

\section{Strong passwords and MFA}
A \textbf{strong password} is long and mixes uppercase letters, lowercase
letters, numbers, and symbols. It must be unique per site, never shared, and
never left at its factory default: routers and office devices ship with
well-known default passwords that must be changed immediately. A
\textbf{password manager} stores one distinct strong password per site so the
user only memorises a single master secret.

\textbf{MFA} stands for Multi-Factor Authentication. It demands two or more
independent proofs of identity: something you know (the password), plus
something you have (a one-time code on the phone) or something you are (a
fingerprint). Even when the password leaks, the attacker still lacks the
second factor, so the account holds. Note the distinction the exam tests: a
longer password is still one factor; MFA is the extra check, not a stronger
password.

\begin{definitionbox}{MFA (Multi-Factor Authentication)}
An authentication method that requires two or more independent factors ---
for example, a password plus a phone code or fingerprint.
\end{definitionbox}

\section{Backup: the availability control}
A \textbf{backup} is a separate clean copy of important data, kept so work
can continue when the original is lost, damaged, or locked by ransomware.
Backup is the great protector of availability. Good backups are regular (so
little work is lost), kept offline or off-site (so one fire, theft, or
ransomware attack cannot destroy both copies), and tested by trial restores. A
backup that cannot be restored is not a backup.

The backup--antivirus boundary is another tested pair. Backup restores
\emph{lost} data; antivirus removes \emph{malicious} data. A machine can have
fresh antivirus and still lose a year's work to a failed disk if there is no
backup.

\section{Safe browsing}
Most infections arrive through the browser, so browsing habits are a control
in their own right. Look for \textbf{HTTPS} and the padlock symbol before
entering passwords or bank details; HTTPS encrypts the session between the
browser and the site. Treat unexpected links and attachments as hostile:
\textbf{phishing} messages imitate trusted senders (a bank, an office, a
colleague) to steal credentials. Download software only from official or
trusted sources, log out on shared computers, and never save passwords on a
machine others use.

\section{Secure Wi-Fi}
The wireless link needs its own protection. Home and office networks must use
\textbf{WPA2 or WPA3} encryption with a strong password, and the router's
default admin password must be changed. Open (password-free) public hotspots
expose traffic, so sensitive logins --- banking, email, office accounts ---
should wait for mobile data, a trusted network, or a \textbf{VPN} (Virtual
Private Network), which encrypts traffic across untrusted links. Two extra
habits help: turn off automatic connection to unknown hotspots, and keep
software updated so known wireless flaws are patched. Note the layering the
exam tests: HTTPS protects the website session, while WPA2/WPA3 protects the
wireless link itself; both layers matter.

\begin{definitionbox}{VPN (Virtual Private Network)}
An encrypted tunnel that carries traffic safely across an untrusted network.
\end{definitionbox}

\section{Removable-media safety}
USB drives and memory cards carry malware between machines. Scan every drive
with antivirus \emph{before} opening its files, and never plug an unknown
found drive into the office computer. Disable \textbf{autorun} so a
newly connected drive cannot launch programs by itself. Handle drives gently
at the end as well: use the safe-eject (safely remove hardware) option before
pulling the device out, so files are not corrupted.

\section{Which control does what}
This table is the highest-yield page in the lesson; most questions are
matching items built from its rows.

\begin{center}
\begin{tabular}{p{0.32\textwidth}p{0.58\textwidth}}
\toprule
\textbf{Control} & \textbf{Main job} \\
\midrule
Antivirus + signatures & Detect and remove known malware files. \\
Real-time scanning + quarantine & Check on access; isolate suspects before harm. \\
Firewall & Block unauthorised network traffic. \\
Patching / updates & Close known software weaknesses. \\
Strong password + MFA & Prove identity; block account theft. \\
Backup & Restore data after loss or ransomware. \\
HTTPS + safe browsing & Protect sessions; dodge phishing and bad downloads. \\
WPA2/WPA3 + VPN & Protect data travelling on the network link. \\
Scan USB + disable autorun & Stop malware entering on removable drives. \\
\bottomrule
\end{tabular}
\end{center}

A useful habit: read the stem's verb. ``Remove'' points to antivirus,
``block traffic'' to the firewall, ``prove identity'' to passwords and MFA,
``restore'' to backup, and ``protect the link'' to WPA2/WPA3 or a VPN.

\section{Exam retrieval and common confusions}
Six distinctions prevent most errors on this topic.
\begin{enumerate}
  \item The firewall blocks \textbf{traffic}; the antivirus removes
    \textbf{files}. Do not swap them.
  \item Quarantine \textbf{isolates}; deletion \textbf{removes}. A
    quarantined file can be restored; a deleted one cannot.
  \item MFA adds a \textbf{second proof of identity}; a longer password is
    still one factor.
  \item Backup restores \textbf{availability}; encryption protects
    \textbf{confidentiality}. A backup is not an antivirus and an antivirus
    is not a backup.
  \item HTTPS protects the \textbf{website session}; WPA2/WPA3 protects the
    \textbf{Wi-Fi link}.
  \item Confidentiality keeps secrets secret, integrity keeps data correct,
    and availability keeps systems reachable. Map each attack to the leg it
    breaks before choosing the control.
\end{enumerate}

\begin{example}
An item asks which control lets the office resume work after ransomware locks
every document. Trace the verb: ``resume work after loss'' means restoring
availability, so the answer is the offline backup. Antivirus (removes malware)
and the firewall (blocks traffic) cannot unlock or recreate the lost files.
\end{example}

\subsection*{Exercises}
\begin{enumerate}[label=\arabic*.]
  \item State the three goals of the CIA triad in one line each.
  \item Explain what a virus signature is and why it must be updated.
  \item Distinguish quarantine from deletion.
  \item Distinguish a firewall from antivirus software.
  \item Explain why a long password alone is not MFA, and what MFA adds.
\end{enumerate}

\subsection*{Solutions}
\begingroup\small
\begin{enumerate}[label=\arabic*.]
  \item Confidentiality: only authorised persons can read the data.
    Integrity: the data is accurate and unaltered by unauthorised means.
    Availability: the data and systems are reachable when needed.
  \item A signature is a stored pattern of a known virus used for recognition.
    New malware appears daily, so the database must be updated or new threats
    will not match anything on record.
  \item Quarantine isolates a suspicious file in a safe area where it cannot
    run, and the file can later be restored; deletion removes the file, and
    it cannot be recovered from quarantine-style holding.
  \item A firewall monitors network traffic and blocks what its rules forbid;
    it does not clean infected files. Antivirus detects and removes malicious
    files on the machine; it does not by itself patrol the network boundary.
  \item A long password is still a single factor (something you know). MFA
    (Multi-Factor Authentication) adds at least one independent factor, such
    as a phone code or fingerprint, so a leaked password alone is not enough.
\end{enumerate}
\endgroup
\end{document}
